% !TeX program = lualatex
% !TeX root = pred4.tex
% ============================================================
% Zajednički preamble za sva predavanja RWA kolegija
% Uključiti s: % ============================================================
% Zajednički preamble za sva predavanja RWA kolegija
% Uključiti s: % ============================================================
% Zajednički preamble za sva predavanja RWA kolegija
% Uključiti s: \input{../common/preamble.tex}
% ============================================================

\documentclass[aspectratio=149,11pt]{beamer}
\usetheme{CambridgeUS}

\usepackage{fontspec}              % radi samo sa XeLaTeX/LuaLaTeX
\setsansfont{Latin Modern Sans}    % sans-serif font
\setmonofont{Latin Modern Mono}    % monospace font za kod
\usepackage[croatian]{babel}
\usepackage{csquotes}              % \enquote{} → hrvatski navodnici „..."
\usepackage{lmodern}
\usepackage{booktabs}
\usepackage{listings}
\usepackage{xcolor}
\usepackage{hyperref}

% ---------- Autor i logo ----------
\author{mag.ing. Josip Torić (SIT UNIZD)}
\titlegraphic{\includegraphics[height=2.5cm]{../common/SIT_logo_tekst1.pdf}}

% ---------- Listings ----------
\definecolor{codebg}{RGB}{248,248,248}
\definecolor{codeframeclr}{RGB}{220,220,220}
\lstdefinelanguage{bash}{
  morekeywords={sudo,apt,brew,pip,pipx,uv,winget,python,git,flask,source,deactivate,code,gunicorn,docker,compose,heroku,render,railway,node,npm,pnpm,yarn,uvicorn,alembic,pytest,ruff,black,precommit},
  sensitive=false,
  morecomment=[l]{\#}
}
\lstdefinelanguage{python}{
  morekeywords={from,import,def,class,return,if,elif,else,for,in,while,try,except,raise,with,as,not,and,or,is,lambda,yield,True,False,None,await,async},
  sensitive=true,
  morecomment=[l]{\#},
  morestring=[b]',
  morestring=[b]"
}
\lstdefinelanguage{sql}{
  morekeywords={SELECT,FROM,WHERE,INSERT,INTO,VALUES,UPDATE,SET,DELETE,CREATE,TABLE,PRIMARY,KEY,FOREIGN,REFERENCES,UNIQUE,NOT,NULL,INDEX,ON,ALTER,DROP,CONSTRAINT},
  sensitive=false,
  morecomment=[l]{--}
}
\lstdefinelanguage{json}{
  morecomment=[l]{//},
  morestring=[b]",
  showstringspaces=false
}
\lstdefinelanguage{typescript}{
  morekeywords={import,export,from,const,let,var,function,return,if,else,for,of,in,while,switch,case,break,async,await,try,catch,finally,throw,new,class,this,extends,implements,interface,type,enum,readonly,public,private,protected,static,get,set,true,false,null,undefined,void,string,number,boolean,any,unknown,never,typeof,keyof,as},
  sensitive=true,
  morecomment=[l]{//},
  morecomment=[s]{/*}{*/},
  morestring=[b]',
  morestring=[b]",
  morestring=[b]`
}
\lstalias{javascript}{typescript}
\lstset{
  backgroundcolor=\color{codebg},
  basicstyle=\ttfamily\footnotesize,
  frame=single,
  rulecolor=\color{codeframeclr},
  frameround=tttt,
  breaklines=true,
  tabsize=2,
  showstringspaces=false,
  upquote=true,
  columns=fullflexible
}

% ---------- Helper macros ----------
\newcommand{\codefile}[1]{\texttt{#1}}
\newcommand{\kbd}[1]{\texttt{#1}}

% ---------- Code frame environment ----------
% Automatski dodaje fragile i shrink=5, koristi se umjesto \begin{frame}
% Primjer: \begin{codeframe}{Naslov} ... \end{codeframe}
% Primjer s opcijama: \begin{codeframe}[shrink=15]{Naslov} ... \end{codeframe}
\newenvironment{codeframe}[2][]
  {\begin{frame}[fragile,shrink=5,#1]{#2}}
  {\end{frame}}

% ============================================================

\documentclass[aspectratio=149,11pt]{beamer}
\usetheme{CambridgeUS}

\usepackage{fontspec}              % radi samo sa XeLaTeX/LuaLaTeX
\setsansfont{Latin Modern Sans}    % sans-serif font
\setmonofont{Latin Modern Mono}    % monospace font za kod
\usepackage[croatian]{babel}
\usepackage{csquotes}              % \enquote{} → hrvatski navodnici „..."
\usepackage{lmodern}
\usepackage{booktabs}
\usepackage{listings}
\usepackage{xcolor}
\usepackage{hyperref}

% ---------- Autor i logo ----------
\author{mag.ing. Josip Torić (SIT UNIZD)}
\titlegraphic{\includegraphics[height=2.5cm]{../common/SIT_logo_tekst1.pdf}}

% ---------- Listings ----------
\definecolor{codebg}{RGB}{248,248,248}
\definecolor{codeframeclr}{RGB}{220,220,220}
\lstdefinelanguage{bash}{
  morekeywords={sudo,apt,brew,pip,pipx,uv,winget,python,git,flask,source,deactivate,code,gunicorn,docker,compose,heroku,render,railway,node,npm,pnpm,yarn,uvicorn,alembic,pytest,ruff,black,precommit},
  sensitive=false,
  morecomment=[l]{\#}
}
\lstdefinelanguage{python}{
  morekeywords={from,import,def,class,return,if,elif,else,for,in,while,try,except,raise,with,as,not,and,or,is,lambda,yield,True,False,None,await,async},
  sensitive=true,
  morecomment=[l]{\#},
  morestring=[b]',
  morestring=[b]"
}
\lstdefinelanguage{sql}{
  morekeywords={SELECT,FROM,WHERE,INSERT,INTO,VALUES,UPDATE,SET,DELETE,CREATE,TABLE,PRIMARY,KEY,FOREIGN,REFERENCES,UNIQUE,NOT,NULL,INDEX,ON,ALTER,DROP,CONSTRAINT},
  sensitive=false,
  morecomment=[l]{--}
}
\lstdefinelanguage{json}{
  morecomment=[l]{//},
  morestring=[b]",
  showstringspaces=false
}
\lstdefinelanguage{typescript}{
  morekeywords={import,export,from,const,let,var,function,return,if,else,for,of,in,while,switch,case,break,async,await,try,catch,finally,throw,new,class,this,extends,implements,interface,type,enum,readonly,public,private,protected,static,get,set,true,false,null,undefined,void,string,number,boolean,any,unknown,never,typeof,keyof,as},
  sensitive=true,
  morecomment=[l]{//},
  morecomment=[s]{/*}{*/},
  morestring=[b]',
  morestring=[b]",
  morestring=[b]`
}
\lstalias{javascript}{typescript}
\lstset{
  backgroundcolor=\color{codebg},
  basicstyle=\ttfamily\footnotesize,
  frame=single,
  rulecolor=\color{codeframeclr},
  frameround=tttt,
  breaklines=true,
  tabsize=2,
  showstringspaces=false,
  upquote=true,
  columns=fullflexible
}

% ---------- Helper macros ----------
\newcommand{\codefile}[1]{\texttt{#1}}
\newcommand{\kbd}[1]{\texttt{#1}}

% ---------- Code frame environment ----------
% Automatski dodaje fragile i shrink=5, koristi se umjesto \begin{frame}
% Primjer: \begin{codeframe}{Naslov} ... \end{codeframe}
% Primjer s opcijama: \begin{codeframe}[shrink=15]{Naslov} ... \end{codeframe}
\newenvironment{codeframe}[2][]
  {\begin{frame}[fragile,shrink=5,#1]{#2}}
  {\end{frame}}

% ============================================================

\documentclass[aspectratio=149,11pt]{beamer}
\usetheme{CambridgeUS}

\usepackage{fontspec}              % radi samo sa XeLaTeX/LuaLaTeX
\setsansfont{Latin Modern Sans}    % sans-serif font
\setmonofont{Latin Modern Mono}    % monospace font za kod
\usepackage[croatian]{babel}
\usepackage{csquotes}              % \enquote{} → hrvatski navodnici „..."
\usepackage{lmodern}
\usepackage{booktabs}
\usepackage{listings}
\usepackage{xcolor}
\usepackage{hyperref}

% ---------- Autor i logo ----------
\author{mag.ing. Josip Torić (SIT UNIZD)}
\titlegraphic{\includegraphics[height=2.5cm]{../common/SIT_logo_tekst1.pdf}}

% ---------- Listings ----------
\definecolor{codebg}{RGB}{248,248,248}
\definecolor{codeframeclr}{RGB}{220,220,220}
\lstdefinelanguage{bash}{
  morekeywords={sudo,apt,brew,pip,pipx,uv,winget,python,git,flask,source,deactivate,code,gunicorn,docker,compose,heroku,render,railway,node,npm,pnpm,yarn,uvicorn,alembic,pytest,ruff,black,precommit},
  sensitive=false,
  morecomment=[l]{\#}
}
\lstdefinelanguage{python}{
  morekeywords={from,import,def,class,return,if,elif,else,for,in,while,try,except,raise,with,as,not,and,or,is,lambda,yield,True,False,None,await,async},
  sensitive=true,
  morecomment=[l]{\#},
  morestring=[b]',
  morestring=[b]"
}
\lstdefinelanguage{sql}{
  morekeywords={SELECT,FROM,WHERE,INSERT,INTO,VALUES,UPDATE,SET,DELETE,CREATE,TABLE,PRIMARY,KEY,FOREIGN,REFERENCES,UNIQUE,NOT,NULL,INDEX,ON,ALTER,DROP,CONSTRAINT},
  sensitive=false,
  morecomment=[l]{--}
}
\lstdefinelanguage{json}{
  morecomment=[l]{//},
  morestring=[b]",
  showstringspaces=false
}
\lstdefinelanguage{typescript}{
  morekeywords={import,export,from,const,let,var,function,return,if,else,for,of,in,while,switch,case,break,async,await,try,catch,finally,throw,new,class,this,extends,implements,interface,type,enum,readonly,public,private,protected,static,get,set,true,false,null,undefined,void,string,number,boolean,any,unknown,never,typeof,keyof,as},
  sensitive=true,
  morecomment=[l]{//},
  morecomment=[s]{/*}{*/},
  morestring=[b]',
  morestring=[b]",
  morestring=[b]`
}
\lstalias{javascript}{typescript}
\lstset{
  backgroundcolor=\color{codebg},
  basicstyle=\ttfamily\footnotesize,
  frame=single,
  rulecolor=\color{codeframeclr},
  frameround=tttt,
  breaklines=true,
  tabsize=2,
  showstringspaces=false,
  upquote=true,
  columns=fullflexible
}

% ---------- Helper macros ----------
\newcommand{\codefile}[1]{\texttt{#1}}
\newcommand{\kbd}[1]{\texttt{#1}}

% ---------- Code frame environment ----------
% Automatski dodaje fragile i shrink=5, koristi se umjesto \begin{frame}
% Primjer: \begin{codeframe}{Naslov} ... \end{codeframe}
% Primjer s opcijama: \begin{codeframe}[shrink=15]{Naslov} ... \end{codeframe}
\newenvironment{codeframe}[2][]
  {\begin{frame}[fragile,shrink=5,#1]{#2}}
  {\end{frame}}


\title[Razvoj web aplikacija]{Razvoj web aplikacija - Predavanje 4 \\ Autorizacija, role i ownership}
\date{26.10.2026.}

\begin{document}

\frame{\titlepage}

% ------------------------------------------------------------
\begin{frame}{Plan predavanja}
\begin{enumerate}
  \item Autentikacija vs autorizacija - razlika i granica
  \item Role: admin i club korisnik
  \item Ownership: korisnik smije mijenjati samo svoje podatke
  \item Dependency-ji za provjeru role i ownership
  \item 401 vs 403 - kad koji status code
  \item Middleware vs dependency pristup
  \item Testovi za autorizaciju
\end{enumerate}

\vspace{0.6em}
Na kraju predavanja: endpointi provjeravaju \textbf{tko si} i \textbf{smiješ li to raditi}.
\end{frame}

% ------------------------------------------------------------
\begin{frame}{Ishodi učenja}
Nakon ovog predavanja student može:
\begin{itemize}
  \item objasniti razliku između autentikacije i autorizacije
  \item implementirati provjeru role u FastAPI dependencyju
  \item implementirati ownership provjeru u service layeru
  \item razlikovati 401 i 403 status code i objasniti kad se koji koristi
  \item napisati test koji provjerava da korisnik ne može pristupiti tuđim podacima
\end{itemize}
\end{frame}

% ------------------------------------------------------------
\begin{frame}{Modeli autorizacije: RBAC, ABAC, ACL}
\begin{tabular}{@{}lp{8cm}@{}}
  \toprule
  \textbf{Model} & \textbf{Kako radi} \\
  \midrule
  \textbf{ACL} (Access Control List) & Svaki resurs ima listu tko smije što (npr. Unix file permissions) \\
  \midrule
  \textbf{RBAC} (Role-Based) & Korisnik ima rolu, rola ima dozvole. Jednostavno za male sustave. \\
  \midrule
  \textbf{ABAC} (Attribute-Based) & Odluka na temelju atributa korisnika, resursa i konteksta. Fleksibilno ali složeno. \\
  \bottomrule
\end{tabular}

\vspace{0.6em}
Naš projekt koristi \textbf{RBAC + ownership} - dovoljno za domenu, a studenti razumiju model.
\end{frame}

% ------------------------------------------------------------
\begin{frame}{Zašto RBAC za naš projekt?}
\begin{itemize}
  \item Imamo samo 2 role (\texttt{admin}, \texttt{club}) - RBAC je prirodan izbor
  \item ABAC bi bio overkill: nemamo složene politike poput \enquote{klub smije prijaviti samo liftera starijeg od 18 godina iz istog grada}
  \item ACL ne skalira: ne želimo definirati pristup za svaki resurs pojedinačno
\end{itemize}

\vspace{0.6em}
\begin{block}{Kad bi trebao ABAC?}
Kad imate 10+ rola, politike ovise o vremenu/lokaciji, ili pravilima koja se mijenjaju bez deploya.
\end{block}
\end{frame}

% ------------------------------------------------------------
\begin{frame}{OWASP Top 10: Broken Access Control (\#1)}
Od 2021., \textbf{Broken Access Control} je \#1 ranjivost prema OWASP-u. \\
OWASP je kratica za \textbf{Open Web Application Security Project} - neprofitna organizacija koja objavljuje liste najčešćih sigurnosnih propusta.

\vspace{0.6em}
Tipični propusti:
\begin{enumerate}
  \item Provjera pristupa samo na frontendu
  \item Endpoint bez ikakve provjere role
  \item Korisnik mijenja ID u URL-u i pristupa tuđim podacima
  \item API vraća podatke koje korisnik ne bi smio vidjeti
\end{enumerate}

\vspace{0.6em}
Naš projekt pokriva sve navedeno: backend provjerava rolu i ownership na \textbf{svakom} endpointu.
\end{frame}

% ------------------------------------------------------------
\begin{frame}{IDOR: Insecure Direct Object Reference}
Napadač mijenja ID u requestu da pristupi tuđim podacima.

\vspace{0.6em}
\textbf{Primjer napada:}
\begin{enumerate}
  \item Korisnik kluba Behemot šalje: \texttt{GET /clubs/1/lifters} (svoj klub)
  \item Mijenja URL u: \texttt{GET /clubs/2/lifters} (klub Heraklo)
  \item Bez ownership provjere: vidi tuđe natjecatelje
\end{enumerate}

\vspace{0.6em}
\textbf{Obrana}: usporedba \texttt{current\_user.club\_id} s \texttt{club\_id} iz URL-a u service layeru.

\vspace{0.4em}
IDOR je najčešći oblik Broken Access Controla.
\end{frame}

% ------------------------------------------------------------
\begin{frame}{Horizontalna vs vertikalna eskalacija}
\begin{tabular}{@{}lp{5cm}p{5cm}@{}}
  \toprule
  & \textbf{Horizontalna} & \textbf{Vertikalna} \\
  \midrule
  Što & Pristup podacima \textbf{drugog korisnika} iste role & Pristup funkcijama \textbf{više role} \\
  Primjer & Klub A vidi liftera kluba B & Club korisnik kreira natjecanje (admin akcija) \\
  Obrana & Ownership provjera & Role provjera \\
  \bottomrule
\end{tabular}

\vspace{0.6em}
Moramo se braniti od \textbf{obje} eskalacije. \texttt{require\_role} brani od vertikalne, ownership od horizontalne.
\end{frame}

% ------------------------------------------------------------
\begin{frame}{Principle of Least Privilege}
\begin{center}
\Large Korisnik smije \textbf{samo ono} što mu treba za posao. Ništa više.
\end{center}

\vspace{0.8em}
U praksi:
\begin{itemize}
  \item Club korisnik ne vidi API za kreiranje natjecanja - čak ni u Swagger dokumentaciji
  \item Admin ne modificira prijave klubova direktno - to je klubova odgovornost
  \item Endpointi koji vraćaju liste filtriraju prema roli: admin vidi sve, club vidi svoje
\end{itemize}
\end{frame}

% ------------------------------------------------------------
\begin{frame}{Defense in Depth: višeslojni pristup}
Ne oslanjamo se na jednu provjeru:

\vspace{0.6em}
\begin{enumerate}
  \item \textbf{Transport}: HTTPS (enkriptirani promet)
  \item \textbf{Autentikacija}: JWT token verifikacija
  \item \textbf{Autorizacija - rola}: \texttt{require\_role} dependency
  \item \textbf{Autorizacija - ownership}: service layer provjera
  \item \textbf{Baza}: foreign key i unique constrainti
\end{enumerate}

\vspace{0.6em}
Ako jedan sloj zakaže, problem zaustavlja sljedeći.
\end{frame}

% ------------------------------------------------------------
\begin{frame}{JWT claims: što stavljamo u token?}
Naš JWT payload sadrži:
\begin{itemize}
  \item \texttt{sub} - user ID
  \item \texttt{role} - \texttt{"admin"} ili \texttt{"club"}
  \item \texttt{club\_id} - ID kluba (null za admina)
  \item \texttt{exp} - istek tokena
\end{itemize}

\vspace{0.6em}
\textbf{Pravilo}: u token stavimo samo ono što nam treba za autorizaciju na \textbf{svakom} requestu.

\vspace{0.4em}
\begin{block}{Zašto ne čitamo rolu iz baze na svakom requestu?}
Performanse: upit u bazu na svaki request je skup. JWT je self-contained - sve što treba je u tokenu.
Trade-off: ako se rola promijeni, stari token i dalje ima staru rolu do isteka.
\end{block}
\end{frame}

% ------------------------------------------------------------
\begin{frame}{Token revokacija: trade-off}
Problem: korisnik je deaktiviran, ali ima valjan token do isteka.

\vspace{0.6em}
Opcije:
\begin{enumerate}
  \item \textbf{Kratki tokeni} (15 min) - kompromis brzine i sigurnosti
  \item \textbf{Blacklist} - lista revociranih tokena (zahtijeva storage, npr. Redis)
  \item \textbf{Verzioniranje} - \texttt{token\_version} u bazi, odbijaj stare verzije
\end{enumerate}

\vspace{0.6em}
Naš projekt koristi kratke tokene (15 min access + refresh). Na obrani morate znati objasniti zašto nismo implementirali blacklist.
\end{frame}

% ------------------------------------------------------------
\begin{frame}{Auto-login: kreiranje kluba s korisnikom}
Kad admin kreira klub:
\begin{enumerate}
  \item Admin šalje: \texttt{POST /clubs} s \texttt{name}, \texttt{city}, \texttt{password}
  \item Backend kreira Club i automatski kreira User:
  \begin{itemize}
    \item \texttt{username = slugify(name)} - npr. \enquote{Behemot} $\rightarrow$ \texttt{behemot}
    \item \texttt{role = "club"}
    \item \texttt{club\_id = novi\_klub.id}
  \end{itemize}
  \item Response vraća klub podatke + generirani username
\end{enumerate}

\vspace{0.6em}
Club 1:1 User - svaki klub ima točno jednog korisnika za prijavu.

\vspace{0.4em}
Admin vidi username i prosljeđuje ga klubu. Klub se prijavi s username + password.
\end{frame}

% ------------------------------------------------------------
\begin{frame}{Kontekst: gdje smo u projektu?}
\begin{itemize}
  \item Predavanje 3: korisnik se može \textbf{prijaviti} i dobiti JWT token
  \item \texttt{get\_current\_user} dependency izvlači korisnika iz tokena
  \item Danas: dodajemo sloj koji provjerava \textbf{smiješ li raditi to što tražiš}
\end{itemize}

\vspace{0.6em}
\begin{center}
\Large Autentikacija = tko si ti. \\ Autorizacija = smiješ li to.
\end{center}
\end{frame}

% ------------------------------------------------------------
\begin{frame}{Naš model autorizacije}
\begin{tabular}{@{}lp{9cm}@{}}
  \toprule
  \textbf{Rola} & \textbf{Ovlasti} \\
  \midrule
  \texttt{admin} & Kreira natjecanja, kreira klubove i korisnike, vidi sve prijave, upravlja rokovima \\
  \midrule
  \texttt{club} & Registrira natjecatelje svog kluba, prijavljuje ih na natjecanje, vidi samo \textbf{svoje} podatke \\
  \bottomrule
\end{tabular}

\vspace{0.6em}
Rola je pohranjena u JWT payloadu i u bazi (\texttt{users.role}).
\end{frame}

% ------------------------------------------------------------
\begin{frame}{Ownership: zašto rola nije dovoljna?}
\begin{itemize}
  \item Dva korisnika imaju rolu \texttt{club}
  \item Korisnik kluba \enquote{Behemot} ne smije uređivati natjecatelje kluba \enquote{Heraklo}
  \item Rola kaže \textbf{što smiješ raditi} - ownership kaže \textbf{nad čijim podacima}
\end{itemize}

\vspace{0.6em}
Ownership se provjerava usporedbom \texttt{current\_user.club\_id} s \texttt{club\_id} resursa.
\end{frame}

% ------------------------------------------------------------
\begin{frame}{401 vs 403: kad koji?}
\begin{tabular}{@{}lp{8cm}@{}}
  \toprule
  \textbf{Status} & \textbf{Značenje} \\
  \midrule
  \textbf{401 Unauthorized} & Nisi se identificirao ili ti je token istekao. Ponovno se prijavi. \\
  \midrule
  \textbf{403 Forbidden} & Znam tko si, ali \textbf{nemaš pravo} na to. Ponovna prijava neće pomoći. \\
  \bottomrule
\end{tabular}

\vspace{0.6em}
Česta greška: vratiti 401 kad bi trebao biti 403. To zbunjuje frontend koji pokušava refreshati token umjesto da prikaže \enquote{nemate pristup}.
\end{frame}

% ------------------------------------------------------------
\begin{frame}[fragile,shrink=5]{Dependency: \texttt{require\_role}}
\begin{lstlisting}[language=python]
from app.core.errors import AppError

def require_role(*allowed_roles: str):
    def checker(
        current_user: User = Depends(get_current_user)
    ) -> User:
        if current_user.role not in allowed_roles:
            raise AppError(
                "forbidden",
                "Nemate dozvolu za ovu akciju",
                status_code=403,
            )
        return current_user
    return checker
\end{lstlisting}

\vspace{0.4em}
Factory pattern: \texttt{require\_role("admin")} vraća dependency koji propušta samo admine.
\end{frame}

% ------------------------------------------------------------
\begin{frame}[fragile,shrink=5]{Korištenje \texttt{require\_role} u routeru}
\begin{lstlisting}[language=python]
@router.post("/")
async def create_club(
    body: ClubCreate,
    db: AsyncSession = Depends(get_db),
    admin: User = Depends(require_role("admin")),
):
    # auto-kreira User s username=slugify(name)
    return await club_service.create_club(db, body)

# preview: ovo implementiramo na P5
@router.get("/{club_id}/lifters")
async def list_lifters(
    club_id: int,
    db: AsyncSession = Depends(get_db),
    user: User = Depends(require_role("admin", "club")),
):
    return await lifter_service.list_by_club(db, club_id, user)
\end{lstlisting}
\end{frame}

% ------------------------------------------------------------
\begin{frame}{Gdje provjeravati ownership?}
\begin{tabular}{@{}lp{8cm}@{}}
  \toprule
  \textbf{Sloj} & \textbf{Odgovornost} \\
  \midrule
  Router & Parsira request, poziva service. \textbf{Ne provjerava} ownership. \\
  \midrule
  Service & \textbf{Provjerava ownership} i poslovna pravila. Baca \texttt{AppError}. \\
  \midrule
  Repository & Radi DB upit. Ne zna za role ni ownership. \\
  \bottomrule
\end{tabular}

\vspace{0.6em}
Ownership je \textbf{poslovno pravilo}, ne HTTP koncept. Zato živi u service layeru.
\end{frame}

% ------------------------------------------------------------
\begin{frame}[fragile, shrink=10]{Service: ownership provjera (preview za P5)}
\begin{lstlisting}[language=python]
async def list_lifters_by_club(
    db: AsyncSession,
    club_id: int,
    current_user: User,
) -> list[Lifter]:
    if current_user.role == "club" and current_user.club_id != club_id:
        raise AppError(
            "forbidden",
            "Ne mozete pristupiti natjecateljima drugog kluba",
            status_code=403,
        )
    return await lifter_repo.get_by_club(db, club_id)
\end{lstlisting}

\vspace{0.4em}
Admin smije vidjeti sve. Club korisnik smije vidjeti samo svoj klub.
Lifter endpoint implementiramo na P5, ovdje pokazujemo \textbf{pattern}.
\end{frame}

% ------------------------------------------------------------
\begin{frame}[fragile,shrink=5]{Alternativa: ownership kao dependency}
\begin{lstlisting}[language=python]
async def require_club_access(
    club_id: int,
    current_user: User = Depends(get_current_user),
) -> User:
    if current_user.role == "admin":
        return current_user
    if current_user.club_id != club_id:
        raise AppError("forbidden", "Pristup odbijen", 403)
    return current_user
\end{lstlisting}

\vspace{0.4em}
Ovo radi ako je \texttt{club\_id} uvijek u URL-u. Za složenije slučajeve ownership ostaje u serviceu.
\end{frame}

% ------------------------------------------------------------
\begin{frame}{Admin vidi sve, club vidi svoje}
Tipični pattern u našem projektu (neki endpointi dolaze na P5):

\vspace{0.6em}
\begin{tabular}{@{}llll@{}}
  \toprule
  \textbf{Endpoint} & \textbf{Admin} & \textbf{Club} & \textbf{Provjera} \\
  \midrule
  \texttt{GET /clubs/\{id\}/lifters} & svi & samo svoj klub & ownership \\
  \texttt{POST /competitions/\{id\}/registrations} & sve & samo svoj klub & ownership \\
  \texttt{POST /competitions} & da & ne & role \\
  \texttt{GET /clubs} & svi & samo svoj & ownership \\
  \texttt{POST /clubs} & da & ne & role \\
  \bottomrule
\end{tabular}
\end{frame}

% ------------------------------------------------------------
\begin{frame}{Česte greške u autorizaciji}
\begin{enumerate}
  \item Provjera role samo na frontendu - backend je autoritet
  \item Ownership provjera samo po URL-u, ne po tijelu requesta
  \item Zaboraviti provjeriti ownership na UPDATE/DELETE
  \item Vratiti 404 umjesto 403 - otkriva postojanje resursa
  \item Ne testirati autorizaciju uopće
\end{enumerate}

\vspace{0.6em}
Na obrani testiram točno ovo: šaljem request s tokenom korisnika koji \textbf{ne bi smio} imati pristup.
\end{frame}

% ------------------------------------------------------------
\begin{frame}{404 vs 403: dilema otkrivanja resursa}
\begin{itemize}
  \item Ako vratim 403: napadač zna da resurs postoji
  \item Ako vratim 404: napadač ne zna, ali korisnik ne razumije zašto ne vidi podatke
\end{itemize}

\vspace{0.6em}
Uobičajena praksa:
\begin{itemize}
  \item Za javne resurse (natjecanja): 403 je ok, ionako su javni
  \item Za osjetljive resurse: 404 je sigurnija opcija
\end{itemize}

\vspace{0.4em}
U našem projektu koristimo 403 jer domena nije osjetljiva. Na obrani morate znati objasniti zašto.
\end{frame}

% ------------------------------------------------------------
\begin{frame}[fragile, shrink=10]{Test: autorizacija po roli}
\begin{lstlisting}[language=python]
def test_club_user_cannot_create_club(client, club_token):
    r = client.post(
        "/clubs",
        json={"name": "Novi", "city": "Zagreb", "password": "pass123"},
        headers={"Authorization": f"Bearer {club_token}"},
    )
    assert r.status_code == 403

def test_admin_can_create_club(client, admin_token):
    r = client.post(
        "/clubs",
        json={"name": "Novi", "city": "Zagreb", "password": "pass123"},
        headers={"Authorization": f"Bearer {admin_token}"},
    )
    assert r.status_code == 201
    assert r.json()["username"] == "novi"
\end{lstlisting}
\end{frame}

% ------------------------------------------------------------
\begin{frame}[fragile, shrink=10]{Test: ownership}
\begin{lstlisting}[language=python]
def test_club_cannot_see_other_clubs_lifters(
    client, club_a_token, club_b_id
):
    r = client.get(
        f"/clubs/{club_b_id}/lifters",
        headers={"Authorization": f"Bearer {club_a_token}"},
    )
    assert r.status_code == 403

def test_club_can_see_own_lifters(
    client, club_a_token, club_a_id
):
    r = client.get(
        f"/clubs/{club_a_id}/lifters",
        headers={"Authorization": f"Bearer {club_a_token}"},
    )
    assert r.status_code == 200
\end{lstlisting}
\end{frame}

% ------------------------------------------------------------
\begin{frame}{Testne fixture za auth tokene}
Testovi trebaju tokene za razne korisnike. Fixture pristup:
\begin{itemize}
  \item \texttt{admin\_token} - fixture koja seedira admina i vrati token
  \item \texttt{club\_a\_token} - fixture za korisnika kluba A
  \item \texttt{club\_b\_token} - fixture za korisnika kluba B
\end{itemize}

\vspace{0.6em}
Alternativa: \texttt{dependency\_overrides} iz prošlog predavanja.
Direktni tokeni su bolji jer testiraju \textbf{cijeli stack} uključujući JWT verifikaciju.
\end{frame}

% ------------------------------------------------------------
\begin{frame}{Middleware vs dependency za autorizaciju}
\begin{tabular}{@{}lp{5cm}p{5cm}@{}}
  \toprule
  & \textbf{Middleware} & \textbf{Dependency} \\
  \midrule
  Kad se izvršava & Prije routera, na svaki request & Samo na endpointima koji ga koriste \\
  Granularnost & Gruba (po prefixu URL-a) & Fina (po endpointu) \\
  Pristup podacima & Samo request/response & Može primiti path parametre \\
  \bottomrule
\end{tabular}

\vspace{0.6em}
FastAPI dependency pristup je bolji za naš slučaj:
\begin{itemize}
  \item Različiti endpointi trebaju različite role
  \item \texttt{club\_id} iz URL-a nam treba za ownership provjeru
  \item Middleware ne može znati koji je \texttt{club\_id} relevantan
\end{itemize}
\end{frame}

% ------------------------------------------------------------
\begin{frame}{Audit checklist: jeste li sve pokrili?}
Prije obrane provjerite svaki endpoint:
\begin{enumerate}
  \item Zahtijeva li autentikaciju? (\texttt{get\_current\_user} ili \texttt{require\_role})
  \item Provjerava li rolu? (admin-only vs club-allowed)
  \item Provjerava li ownership? (club smije samo svoje)
  \item Postoji li test za neautorizirani pristup?
  \item Je li response konzistentan? (401 vs 403 na pravom mjestu)
\end{enumerate}

\vspace{0.4em}
Ako imate 5 endpointa, trebate barem 5 testova za autorizaciju - po jedan negativni za svaki endpoint.
\end{frame}

% ------------------------------------------------------------
\begin{frame}{Commitovi za Predavanje 4}
\begin{itemize}
  \item \texttt{feat(api): add require\_role dependency factory}
  \item \texttt{feat(api): add club schemas and repository}
  \item \texttt{feat(api): add admin club endpoints (create with auto-login, list, update, reset-password)}
  \item \texttt{feat(api): enforce ownership checks in service layer}
  \item \texttt{test(api): add role and ownership tests (admin/club/cross-club)}
\end{itemize}

\vspace{0.6em}
\textbf{Na satu:} prolazimo diffeve i testiramo sa Swaggerom - šaljemo requestove s raznim tokenima.
\end{frame}

% ------------------------------------------------------------
\begin{frame}{Što ide sljedeće (Predavanje 5)}
\begin{itemize}
  \item Modeliranje preostalih entiteta: Lifter, Competition, Registration
  \item CRUD operacije kroz sve slojeve
  \item Pydantic schema dizajn: Create vs Update vs Response
  \item Validacija inputa i smisleni error odgovori
  \item Paginacija i filtriranje
\end{itemize}

\vspace{0.8em}
\begin{center}
\Large Pitanja?
\end{center}
\end{frame}

\end{document}
